% Part: model-theory
% Chapter: interpolation
% Section: definability

\documentclass[../../../include/open-logic-section]{subfiles}

\begin{document}

\olfileid[tr]{mod}{int}{def}

\olsection{Tanımlanabilirlik Teoremi}

Aradeğerleme teoreminin önemli bir uygulaması Beth tanımlanabilirlik
teoremidir. $n$ yerli bir $R$ bağıntısını tanımlamak için $R$'yi
içermeyen ve $n$ serbest !!{variable}s içeren !!a{formula}~$!C$
verebiliriz. Bu, $R$'nin $!C$ cinsinden \emph{açık} bir tanımı olur.
Ardından, $n$ yerli !!{predicate}~$P$ içeren !!a{language} içindeki
bir $\Sigma(P)$ kuramının $P$'yi açıkça tanımladığını, biçimselleştirilmiş
bir açık tanımı içeriyorsa (ya da en azından gerektiriyorsa)
söyleyebiliriz; yani
\[
\Sigma(P) \Entails \lforall[x_1][\dots
  \lforall[x_n][(\Atom{P}{x_1,\dots, x_n} \liff !C(x_1, \dots,
    x_n))]].
\]
Fakat açık tanım, bir bağıntıyı tam olarak belirleme anlamında
tanımlamanın yalnızca bir yoludur. Bir kuram, her modelde $P$'nin
yorumunun !!{language} dilinin geri kalanının yorumuyla sabitlendiği
bir kuram da olabilir. Tanımlanabilirlik teoremi, bir kuram $P$'nin
yorumunu bu biçimde sabitlediğinde, yani $P$'yi \emph{örtük olarak
tanımladığında}, onu açıkça da tanımladığını söyler.

\begin{defn}
$\Lang L$, !!{predicate}~$P$ içermeyen !!a{language} olsun.
$\Lang L\cup\{P\}$ dilindeki bir !!{sentence}s kümesi~$\Sigma(P)$,
ancak ve ancak $\Lang L$ dilinde
\[
\Sigma(P) \Entails \lforall[x_1][\dots
  \lforall[x_n][(\Atom{P}{x_1,\dots, x_n} \liff !C(x_1, \dots,
    x_n))]]
\]
olacak !!a{formula}~$!C(x_1,\dots,x_n)$ varsa $P$'yi
\emph{açıkça tanımlar}.
\end{defn}

\begin{defn}
$\Lang L$, !!{predicate}s~$P$ ve $P'$ simgelerini içermeyen
!!a{language} olsun. $\Lang L\cup\{P\}$ dilindeki bir !!{sentence}s
kümesi~$\Sigma(P)$, ancak ve ancak
\[
\Sigma(P) \cup \Sigma(P') \Entails \lforall[x_1][\dots
  \lforall[x_n][(\Atom{P}{x_1,\dots, x_n} \liff \Atom{P'}{x_1,\dots,
      x_n})]]
\]
ise $P$'yi \emph{örtük olarak tanımlar}. Burada $\Sigma(P')$,
$\Sigma(P)$ içinde $P$'nin her yerde $P'$ ile değiştirilmesiyle elde
edilir.
\end{defn}

Başka bir deyişle her $\Struct M$ modeli ve
$R,R'\subseteq\Domain M^n$ için hem
$\Sat{\Expan{M}{R}}{\Sigma(P)}$ hem
$\Sat{\Expan{M}{R'}}{\Sigma(P')}$ ise $R=R'$ olur. Burada
$\Expan M R$, $\Lang L$ dilinin $\Lang L\cup\{P\}$ diline,
$\Assign P{M'}=R$ olacak biçimde genişlemesi olan
!!{structure}~$\Struct{M'}$; $\Expan M{R'}$ da benzeridir.

\begin{thm}[Beth tanımlanabilirlik teoremi]
$\Lang L\cup\{P\}$ dilindeki bir $\Sigma(P)$ !!{formula}s kümesi
$P$'yi örtük olarak tanımlar ancak ve ancak açıkça tanımlar.
\end{thm}

\begin{proof}
$\Sigma(P)$, $P$'yi açıkça tanımlıyorsa
\begin{align*}
  \Sigma(P) & \Entails & \lforall[x_1][\dots \lforall[x_n]
    [(\Atom{P}{x_1,\dots, x_n} \liff !C(x_1,\dots,x_n))]]\\
  \Sigma(P') & \Entails & \lforall[x_1][\dots \lforall[x_n]
    [(\Atom{P'}{x_1,\dots, x_n} \liff !C(x_1,\dots,x_n))]]
\end{align*}
olur ve sonuç çıkar. Ters yön için $\Sigma(P)$'nin $P$'yi örtük olarak
tanımladığını varsayalım. Önce $\Lang L$ diline !!{constant}s
$c_1,\dots,c_n$ ekleyelim. O zaman
\[
\Sigma(P) \cup \Sigma(P') \Entails
\Atom{P}{c_1, \dots, c_n} \to \Atom{P'}{c_1, \dots, c_n}.
\]
Tıkızlık gereği sonlu $\Delta_0\subseteq\Sigma(P)$ ve
$\Delta_1\subseteq\Sigma(P')$ kümeleri vardır ve
\[
\Delta_0 \cup \Delta_1 \Entails
\Atom{P}{c_1, \dots, c_n} \to \Atom{P'}{c_1, \dots, c_n}.
\]
$!D(P)$, ya $!A(P)\in\Delta_0$ ya da $!A(P')\in\Delta_1$ olan bütün
!!{sentence}s~$!A(P)$ ifadelerinin birleşimi; $!D(P')$ de aynı
koşuldaki bütün !!{sentence}s~$!A(P')$ ifadelerinin birleşimi olsun.
O zaman
$!D(P)\land !D(P')\Entails
\Atom P{c_1,\dots,c_n}\to P'c_1\dots c_n$. Her !!{predicate}
tek bir $\Entails$ tarafında bulunacak biçimde bunu yeniden düzenleriz:
\[
!D(P) \land \Atom{P}{c_1, \dots, c_n} \Entails
!D(P') \to \Atom{P'}{c_1, \dots, c_n}.
\]
Craig aradeğerleme teoremi gereği $P$ veya $P'$ içermeyen
!!a{sentence}~$!C(c_1,\dots,c_n)$ vardır ve
\begin{align*}
  !D(P) \land \Atom{P}{c_1, \dots, c_n}
    & \Entails !C(c_1,\dots, c_n); \\
  !C(c_1,\dots, c_n)
    & \Entails !D(P') \to \Atom{P'}{c_1, \dots, c_n}).
\end{align*}
Birinci gerektirmeden
$!D(P)\Entails\Atom P{c_1,\dots,c_n}\lif !C(c_1,\dots,c_n)$ elde
edilir. İkinciden de, bir $\Lang L\cup\{P\}$ modeli
$\Expan M R$ içinde $!A(P)$ doğruysa ve ancak ona karşılık gelen
$\Lang L\cup\{P'\}$ modelinde $!A(P')$ doğruysa,
$!C(c_1,\dots,c_n)\Entails
!D(P)\lif\Atom{P}{c_1,\dots,c_n}$; buradan
\[
!D(P) \Entails !C(c_1,\dots,c_n) \to \Atom{P}{c_1,\dots, c_n}.
\]
İkisini birleştirince
$!D(P)\Entails\Atom P{c_1,\dots,c_n}\liff !C(c_1,\dots,c_n)$; tekdüzelik
ve genelleme yoluyla da
\[
\Sigma(P) \Entails
\lforall[x_1][\dots\lforall[x_n][(\Atom{P}{x_1,\dots, x_n} \liff
    !C(x_1,\dots, x_n))]]
\]
elde edilir.
\end{proof}

\end{document}
